\section{Introduction}

Automated visual inspection is essential for weld-quality control in high-throughput manufacturing. Annular welds are especially challenging because the seam is a closed, periodically repeating structure, whereas reflectance, width, and illumination vary around the circumference. Operationally important defects can be rare, weak, or incomplete. A detector may consequently assign a low score to a true tail defect for the same visual reason that it assigns a low score to a normal local variation. Once a fixed threshold removes the candidate, downstream classification cannot recover it.

Long-tailed detection has mainly addressed imbalance through training schedules, additional supervision, or classifier adjustment. SimLTD transfers head-class knowledge to tail classes through staged training~\cite{tran2025simltd}. RichSem introduces image-level semantics and coarse locations as supplementary supervision~\cite{meng2023richsem}. LOS revisits classifier re-training through label over-smoothing~\cite{sun2025los}. These methods address important sources of class bias. In industrial inspection, however, a prior query-level decision remains unresolved: whether a low-score candidate contains defect evidence or reflects a normal seam variation. Lowering the score threshold increases candidate coverage, but it also admits reflections, texture fluctuations, and tail-to-head confusions. The resulting problem is one of candidate-evidence reliability rather than score calibration alone.

Spatial and frequency representations offer complementary observations, yet conventional addition or concatenation assigns them fixed relative importance. That assumption is unsuitable for metallic welds. A reflection can create a large luminance deviation without a stable spectral transition. A small pore can be more evident at a high-frequency boundary than in colour, whereas a smoke-like defect can be dominated by gradual low- or middle-frequency variation. Moreover, absolute responses depend on circumferential position because normal seam texture changes around the ring. The central problem is therefore not simply to add a frequency branch. It is to determine which spatial or frequency evidence is reliable for each query after comparison with a compatible normal seam segment.

We address this problem with \emph{NORA}, illustrated in Fig.~\ref{fig:cyclicibo}. NORA treats a D-FINE-S query~\cite{peng2025dfine} as a candidate hypothesis supported by two conditionally reliable evidence streams. A geometry-preserving front end unwraps the annular seam without radial stretching and forms periodic 150-pixel windows. For each query, inside--boundary--outside (IBO) regions describe how colour, gradients, deep spatial features, and local multi-band responses change from the candidate interior to its boundary and neighbouring seam. The transitions are normalized with normal-reference statistics from comparable circumferential locations. A reliability gate then combines the two domains using branch uncertainty and reference-normalized anomaly strength. An explicit disagreement component retains informative cross-domain conflicts. Finally, same-class cyclic merging and restricted correction are applied only to recurrent tail-to-head confusions identified by an error audit.

The contributions are threefold:
\begin{itemize}
    \item We formulate long-tailed annular weld inspection as an evidence-reliability problem. The central method is a query-level spatial--frequency fusion rule whose weights depend on branch uncertainty and normal-reference anomaly strength, while spatial--frequency disagreement is retained as discriminative evidence.
    \item We introduce normal-reference IBO transitions in both domains. The representation measures interior-to-boundary-to-seam changes against comparable normal locations rather than relying on absolute brightness, texture, or whole-image high-frequency energy.
    \item We combine geometry-preserving cyclic candidate generation with same-class cyclic merging and restricted confusion-pair correction. We also report a long-tail audit that separates correct, wrong-class, and missed instances.
\end{itemize}
